\paragraph{倍映射在 \texorpdfstring{$t$}{t}-线上诱导的 Latt\`es 自映射：闭式、临界分解与 Möbius 对称}
在判别式椭圆曲线 \(E_0:y^2=f(t)\) 的 \(t\)-坐标中，倍映射 \([2]\) 诱导一个次数 \(4\) 的有理自映射。
其后临界集由 \(2\)-挠点的 \(t\)-坐标完全控制，并且 \(E_0(\QQ)[2]\) 的平移在 \(t\)-线上实现为 Klein 四元反演族。

\begin{theorem}[倍映射诱导的 \texorpdfstring{$t$}{t}-线 Latt\`es 映射的闭式与后临界集]\label{thm:xi-j-sextic-elliptic-lattes-t-doubling-pcf}
取定理 \ref{thm:xi-j-sextic-discriminant-quadratic-subfield-ellipticization} 中的椭圆曲线
\[
E_0:\quad y^2=(t-1728)\bigl(t^2+1862t+161792\bigr)
=t^3+134t^2-3055744t-279576576.
\]
令 \(\mathcal L:\PP^1_t\dashrightarrow \PP^1_t\) 为由倍映射 \([2]:E_0\to E_0\) 在 \(t\)-坐标下诱导的有理自映射，即对一般点 \(P\in E_0\) 定义
\(
\mathcal L(t(P)):=t([2]P)
\)。
则 \(\deg(\mathcal L)=4\)，并且具有完全显式表达式
\[
\boxed{\;
\mathcal L(t)=
\frac{t^{4}+6111488\,t^{2}+2236612608\,t+9487424438272}
{4\,(t-1728)\,(t^{2}+1862t+161792)}\; }.
\]
其导数满足可约分解
\[
\boxed{\;
\mathcal L'(t)=
\frac{(t^{2}-3456t-3379328)\,
(t^{4}+3724t^{3}+970752t^{2}+10348004864t-8393927966720)}
{4\,(t-1728)^{2}\,(t^{2}+1862t+161792)^{2}}\; }.
\]
从而 \(\mathcal L\) 的有限临界值集合严格等于
\[
\boxed{\ \{1728,\ \alpha,\ \beta\}\ },
\qquad
\alpha,\beta\ \text{为}\ t^{2}+1862t+161792=0\ \text{的两根}.
\]
并且后临界集合为
\[
\boxed{\ \mathrm{PostCrit}(\mathcal L)=\{\infty,\,1728,\,\alpha,\,\beta\}\ },
\]
且所有临界轨道在至多两步内进入不动点 \(\infty\)。
\end{theorem}
\begin{proof}
将 \(E_0\) 写成 Weierstrass 形式 \(y^2=t^3+a_2t^2+a_4t+a_6\)，其中
\(
a_2=134,\ a_4=-3055744,\ a_6=-279576576
\)。
对一般点 \(P=(t,y)\)，切线斜率为
\[
m=\frac{3t^{2}+2a_2t+a_4}{2y}=\frac{3t^{2}+268t-3055744}{2y},
\]
倍点公式给出
\[
t([2]P)=m^{2}-a_2-2t.
\]
将 \(y^2=(t-1728)(t^2+1862t+161792)\) 代入并化简，得到所示 \(\mathcal L(t)\)。
对 \(\mathcal L\) 求导并因式分解得到 \(\mathcal L'(t)\) 的分解式。
\par
二次因子 \(t^{2}-3456t-3379328\) 的任一根 \(\xi\) 满足 \(\xi^{2}=3456\xi+3379328\)，代回 \(\mathcal L(t)\) 可直接化简得 \(\mathcal L(\xi)=1728\)，故 \(1728\) 为临界值。
对四次因子 \(g_4(t)=t^{4}+3724t^{3}+970752t^{2}+10348004864t-8393927966720\)，
以消去恒等式
\[
\mathrm{Res}_{t}\!\bigl(g_4(t),\,\mathrm{num}(\mathcal L)(t)-z\,\mathrm{den}(\mathcal L)(t)\bigr)
=\mathrm{const}\cdot(z^{2}+1862z+161792)^{2}
\]
刻画其像：任意 \(g_4(\eta)=0\) 的根 \(\eta\) 的像必落在 \(\{\alpha,\beta\}\)。
由于 \(\mathcal L\) 在 \(t=1728,\alpha,\beta\) 处有极点，故 \(\mathcal L(1728)=\mathcal L(\alpha)=\mathcal L(\beta)=\infty\)。
又由 \(\deg\mathrm{num}=4\)、\(\deg\mathrm{den}=3\) 得 \(\mathcal L(t)\sim \frac14 t\)（\(t\to\infty\)），从而 \(\mathcal L(\infty)=\infty\)。
综上得到后临界集合与“至多两步入 \(\infty\)”的断言。证毕。
\end{proof}

\begin{theorem}[\(\QQ\)-有理 \(2\)-挠平移诱导的降阶因子化与 \texorpdfstring{$\sqrt{89}$}{sqrt(89)} 临界值]\label{thm:xi-j-sextic-elliptic-lattes-degree-drop-by-2torsion}
仍取 \(E_0\) 与 \(\mathcal L\) 如定理 \ref{thm:xi-j-sextic-elliptic-lattes-t-doubling-pcf}。
记 \(\QQ\)-有理 \(2\)-挠点 \(Q=(1728,0)\in E_0(\QQ)\)，并令 \(\tau_Q:E_0\to E_0\) 为平移 \(P\mapsto P+Q\)。
则存在 involution \(\phi\in\mathrm{PGL}_2(\QQ)\) 使得对一般点 \(P\) 有 \(t(\tau_Q(P))=\phi(t(P))\)，并且
\[
\boxed{\ \phi(t)=1728+\frac{B}{t-1728}\ },
\qquad
B:=(1728-\alpha)(1728-\beta)=6365312.
\]
进一步，倍映射满足强不变性
\[
\boxed{\ \mathcal L\circ \phi=\mathcal L\ }.
\]
令 \(u:=t-1728\) 并定义 \(\phi\)-不变函数
\[
\pi(t):=u+\frac{B}{u},
\qquad (\deg\pi=2),
\]
则存在唯一的次数 \(2\) 有理函数 \(F\) 使得
\[
\boxed{\ \mathcal L = F\circ \pi\ },
\qquad
\boxed{\ F(x)=\frac{x^{2}+6912x+11296768}{4(x+5318)}\ }.
\]
并且
\[
\mathrm{Crit}(F)=\{-5318\pm 178\sqrt{89}\},
\qquad
F\bigl(\mathrm{Crit}(F)\bigr)=\{-931\pm 89\sqrt{89}\}=\{\alpha,\beta\}.
\]
\end{theorem}
\begin{proof}
由于 \(Q\in E_0[2]\)，对任意 \(P\in E_0\) 有
\(
[2](P+Q)=[2]P
\)。
取 \(t\)-坐标即得 \(\mathcal L(t(P+Q))=\mathcal L(t(P))\)，从而 \(\mathcal L\circ\phi=\mathcal L\)。
又 \(E_0\) 的 \(2\)-挠点的 \(t\)-坐标为 \(\{1728,\alpha,\beta\}\)，而 \(t=1728\) 对应的平移在 \(\PP^1_t\) 上为以 \(1728\) 为不动点的分式线性反演；
由对称性与归一化条件（交换 \(\infty\leftrightarrow 1728\) 且保持交比）得到所示 \(\phi\)，其中
\(
B=(1728-\alpha)(1728-\beta)
\)。
\par
\(\phi\) 的不变有理函数域为 \(\QQ(\pi(t))\)，故 \(\mathcal L\) 必经由商映射 \(\pi\) 因子化：存在 \(F\) 使 \(\mathcal L=F\circ\pi\)。
将 \(t=u+1728\) 并代入 \(\pi=u+B/u\)，对 \(\mathcal L(u+1728)\) 作消元化简得到所示 \(F\)；
其导数与判别式给出临界点为 \(-5318\pm 178\sqrt{89}\)，代回 \(F\) 得临界值为 \(-931\pm 89\sqrt{89}\)。
证毕。
\end{proof}

\begin{theorem}[\texorpdfstring{$E_0[2]$}{E0[2]} 的 Klein 平移群在 \texorpdfstring{$t$}{t}-线上实现为 Möbius 反演族]\label{thm:xi-j-sextic-elliptic-lattes-klein-mobius}
令 \(e_1:=1728\)，并令 \(e_2,e_3\) 为二次方程 \(t^{2}+1862t+161792=0\) 的两根（即 \(e_{2,3}=-931\pm 89\sqrt{89}\)）。
对 \(\{i,j,k\}=\{1,2,3\}\) 定义
\[
\phi_i(t):=e_i+\frac{(e_i-e_j)(e_i-e_k)}{t-e_i}.
\]
则：
\begin{enumerate}
  \item \(\phi_i\) 为二阶 Möbius 变换，且在四点集 \(\{\infty,e_1,e_2,e_3\}\) 上的置换为
  \[
  \phi_i:\ (\infty\ e_i)(e_j\ e_k),
  \]
  从而 \(\langle \phi_1,\phi_2,\phi_3\rangle\cong (\ZZ/2\ZZ)^2\) 且 \(\phi_i\circ\phi_j=\phi_k\)（\(i\neq j\)）。
  \item 该 Klein 四元群等同于 \(E_0[2]\) 的平移群在 \(t\)-线上诱导的作用：\(\phi_i\) 对应平移 \(\tau_{(e_i,0)}\)。
  \item 对每个 \(i\in\{1,2,3\}\) 有
  \[
  \boxed{\ \mathcal L\circ \phi_i=\mathcal L\ },
  \]
  因而 \(\mathrm{Aut}_{\overline{\QQ}}(\mathcal L)\) 至少包含该 Klein 四元群；其中在 \(\QQ\) 上可见的 Möbius 对称子群为 \(\langle \phi_1\rangle\)。
\end{enumerate}
\end{theorem}
\begin{proof}
对 (1)，直接代入验证 \(\phi_i(e_i)=\infty\)、\(\phi_i(\infty)=e_i\) 且 \(\phi_i(e_j)=e_k\)。
由此得到 \(\phi_i^2=\mathrm{id}\) 与乘法关系。
\par
对 (2)，在模型 \(y^2=(t-e_1)(t-e_2)(t-e_3)\) 上，\(2\)-挠点平移在 \(t\)-坐标下的作用为上述反演式，亦可由群律加法公式逐项推导。
\par
对 (3)，对任意 \(T\in E_0[2]\) 都有 \([2](P+T)=[2]P\)，取 \(t\)-坐标即得 \(\mathcal L\circ\phi_i=\mathcal L\)。证毕。
\end{proof}

\endinput
